\section{Introducci\'on}

La gesti\'on sostenible de residuos org\'anicos representa un desaf\'io ambiental crucial debido a su contribuci\'on a la emisi\'on de gases de efecto invernadero (GEI), especialmente el di\'oxido de carbono (CO$_2$). En este contexto, la bioconversi\'on de residuos mediante larvas de la mosca soldado negra (\textit{Hermetia illucens}, BSFL) surge como una alternativa eficiente, transformando materia org\'anica en biomasa de alto valor agregado.

Durante el proceso de cr\'ia de BSFL, el CO$_2$ es generado principalmente por la respiraci\'on de las larvas, la actividad metab\'olica de los microorganismos asociados y la biodegradaci\'on del sustrato org\'anico. Parte de estas emisiones de CO$_2$ proviene de fuentes biog\'enicas, tales como la descomposici\'on aer\'obica de la materia org\'anica y la oxidaci\'on de metano (CH$_4$) por bacterias metanotr\'ofas.

El objetivo de este trabajo es desarrollar un modelo matem\'atico que permita predecir la cantidad de CO$_2$ generada a partir de la cantidad inicial de huevos depositados sobre una biomasa disponible. Para ello, se considerar\'an las principales din\'amicas biol\'ogicas del crecimiento larvario, el consumo de sustrato, la producci\'on de CO$_2$ por larvas y microorganismos, y el efecto de variables ambientales como la humedad del sustrato y del ambiente. Esta predicci\'on contribuir\'a a optimizar procesos de manejo de residuos y a reducir las emisiones de GEI en sistemas de producci\'on sostenibles.

\section*{Objetivo General}

Desarrollar y validar un modelo matemático-predictivo para el proceso de bioconversión de residuos orgánicos mediante larvas de \textit{Hermetia illucens}, integrando sensores IoT para medir variables críticas (temperatura, humedad, concentraciones de $CO_2$, $CH_4$ y oxígeno) y algoritmos de aprendizaje automático, con el fin de predecir y optimizar la generación de gases de efecto invernadero, considerando las dinámicas biológicas y la competencia entre larvas y microorganismos, contribuyendo así a la gestión sostenible de residuos y la mitigación de emisiones atmosféricas.
